\chapter{Apêndice B --- Scripts Python Utilizados}
\label{apendiceB}

Este apêndice apresenta trechos representativos dos scripts Python utilizados na
construção do conjunto de dados, no cálculo do Prêmio de Risco de Volatilidade do
Bitcoin (BVRP) e na implementação dos modelos empíricos empregados na dissertação.
O objetivo não é reproduzir integralmente todo o \textit{pipeline} computacional,
mas documentar de forma transparente os principais procedimentos discutidos nos
capítulos anteriores.

Os códigos foram organizados em blocos temáticos: (i) importação de bibliotecas;
(ii) construção das variáveis de volatilidade; (iii) retornos futuros; (iv) BVRP;
(v) classificação de regimes; (vi) regressão por MQO com erros-padrão de Newey--West;
e (vii) estimação dos quatro modelos de aprendizado de máquina com avaliação fora
da amostra. A implementação usa bibliotecas consolidadas em finanças quantitativas e
aprendizado de máquina \citep{andersen2003realized, corsi2009har,
breiman2001randomforest, friedman2001gradientboosting, Gu2020EmpiricalML}.

\section{Importação de bibliotecas}

\begin{verbatim}
import pandas as pd
import numpy as np
import statsmodels.api as sm
from scipy import stats

from sklearn.linear_model import Lasso, Ridge
from sklearn.ensemble import (RandomForestRegressor,
                               GradientBoostingRegressor)
from sklearn.metrics import mean_squared_error, r2_score
from sklearn.preprocessing import StandardScaler

import matplotlib.pyplot as plt
\end{verbatim}

\section{Cálculo da volatilidade realizada (RV30)}

A volatilidade realizada é calculada em janela móvel de 30 dias e anualizada por
365 dias, refletindo a operação contínua (24/7) do mercado de criptoativos:

\begin{verbatim}
def calc_rv30(log_returns):
    """Volatilidade realizada anualizada em janela de 30 dias."""
    rv = np.sqrt(
        365 / 30 * log_returns.rolling(30)
                               .apply(lambda x: np.sum(x**2))
    )
    return rv

df["log_ret"] = np.log(df["close"] / df["close"].shift(1))
df["rv_30d"]  = calc_rv30(df["log_ret"])
\end{verbatim}

\section{Construção dos retornos futuros}

Os retornos futuros são retornos logarítmicos acumulados em cinco horizontes,
com defasagem negativa para evitar \textit{look-ahead bias}:

\begin{verbatim}
for h in [1, 5, 10, 20, 30]:
    df[f"ret_fut_{h}d"] = np.log(
        df["close"].shift(-h) / df["close"]
    )
\end{verbatim}

As variáveis resultantes ($h \in \{1, 5, 10, 20, 30\}$ dias) são utilizadas
nas análises econométricas e preditivas dos
Capítulos~\ref{chap:resultados_ols}, \ref{chap:cap6_ml},
\ref{chap:ml_bvrp} e \ref{chap:regimes}.

\section{Cálculo do BVRP}

O prêmio de risco de volatilidade é a diferença entre volatilidade realizada e
implícita, ambas anualizadas em base de 30 dias (convenção RV $-$ IV,
Decisão~1 --- Fase~0):

\begin{verbatim}
df["vrp_30d"] = df["rv_30d"] - df["iv_30d"]
\end{verbatim}

\section{Classificação de regimes de volatilidade}

Os regimes são definidos pelos percentis 25 e 75 da \texttt{rv\_30d},
calculados sobre a amostra completa:

\begin{verbatim}
q25 = df["rv_30d"].quantile(0.25)   # 39,3% na amostra
q75 = df["rv_30d"].quantile(0.75)   # 65,6% na amostra

df["regime"] = "media"
df.loc[df["rv_30d"] <= q25, "regime"] = "baixa"
df.loc[df["rv_30d"] >= q75, "regime"] = "alta"

# Indicadora binária usada nas regressões condicionais
df["dummy_alta"] = (df["regime"] == "alta").astype(int)
\end{verbatim}

\section{Regressão por MQO com erros-padrão de Newey--West}

As regressões do Capítulo~\ref{chap:resultados_ols} são estimadas com
\texttt{statsmodels}, usando erros-padrão HAC com $h$ defasagens para corrigir
autocorrelação induzida pela sobreposição das janelas de retorno:

\begin{verbatim}
def reg_nw(df, target_col, horizon):
    """OLS com erros-padrão Newey-West (h defasagens)."""
    X = sm.add_constant(df["vrp_30d"])
    y = df[target_col]
    valid = X.join(y).dropna()
    model = sm.OLS(valid[y.name], valid[X.columns])
    return model.fit(cov_type="HAC",
                     cov_kwds={"maxlags": horizon})

for h in [1, 5, 10, 20, 30]:
    res = reg_nw(df, f"ret_fut_{h}d", h)
    print(f"h={h:2d}d  beta={res.params[1]:.4f}  "
          f"t={res.tvalues[1]:.3f}  R2={res.rsquared:.4f}")
\end{verbatim}

Para as regressões condicionais do Capítulo~\ref{chap:regimes}, adiciona-se a
dummy de regime e o termo de interação:

\begin{verbatim}
def reg_condicional(df, h):
    df2 = df.copy()
    df2["bvrp_x_alta"] = df2["vrp_30d"] * df2["dummy_alta"]
    X = sm.add_constant(
        df2[["vrp_30d", "dummy_alta", "bvrp_x_alta"]]
    )
    y = df2[f"ret_fut_{h}d"]
    valid = X.join(y).dropna()
    model = sm.OLS(valid[y.name], valid[X.columns])
    return model.fit(cov_type="HAC",
                     cov_kwds={"maxlags": h})
\end{verbatim}

\section{Modelos de aprendizado de máquina}

\subsection{Preparação e divisão da amostra}

Os quatro modelos são estimados com divisão temporal 70\,\%/30\,\% (treino/teste),
preservando a ordem cronológica. A variável-alvo é o nível do BVRP
(\texttt{vrp\_30d}):

\begin{verbatim}
FEATURES = [
    "ret", "ret_lag_1", "ret_lag_5", "ret_lag_20",
    "d_iv_1d", "d_vrp_1d", "vrp_regime_num",
    "month", "weekday",
]
TARGET = "vrp_30d"

data  = df[FEATURES + [TARGET]].dropna()
split = int(0.7 * len(data))

X_train = data[FEATURES].iloc[:split]
X_test  = data[FEATURES].iloc[split:]
y_train = data[TARGET].iloc[:split]
y_test  = data[TARGET].iloc[split:]
\end{verbatim}

\subsection{LASSO e Ridge}

Os modelos lineares regularizados requerem padronização dos preditores:

\begin{verbatim}
scaler = StandardScaler()
Xtr_s  = scaler.fit_transform(X_train)
Xte_s  = scaler.transform(X_test)

lasso = Lasso(alpha=0.001, max_iter=50000).fit(Xtr_s, y_train)
ridge = Ridge(alpha=1.0).fit(Xtr_s, y_train)

pred_lasso = lasso.predict(Xte_s)
pred_ridge = ridge.predict(Xte_s)
\end{verbatim}

\subsection{Floresta aleatória e Gradient Boosting}

Os modelos baseados em árvores operam diretamente sobre as variáveis não
padronizadas:

\begin{verbatim}
rf = RandomForestRegressor(
    n_estimators=300,
    max_depth=None,
    random_state=42,
).fit(X_train, y_train)

gb = GradientBoostingRegressor(
    n_estimators=300,
    learning_rate=0.05,
    max_depth=4,
    random_state=42,
).fit(X_train, y_train)

pred_rf = rf.predict(X_test)
pred_gb = gb.predict(X_test)
\end{verbatim}

\subsection{Avaliação fora da amostra}

O desempenho é comparado pelo $R^2$ fora da amostra e pela raiz do erro
quadrático médio (RMSE), tomando a média histórica do período de treino como
\textit{benchmark}:

\begin{verbatim}
benchmark = np.full(len(y_test), y_train.mean())

def oos_metrics(y_true, y_pred, label):
    rmse = np.sqrt(mean_squared_error(y_true, y_pred))
    r2   = r2_score(y_true, y_pred)
    print(f"{label:20s}  RMSE={rmse:.4f}  R2={r2:.4f}")

oos_metrics(y_test, benchmark,  "Benchmark (media)")
oos_metrics(y_test, pred_lasso, "LASSO")
oos_metrics(y_test, pred_ridge, "Ridge")
oos_metrics(y_test, pred_rf,    "Random Forest")
oos_metrics(y_test, pred_gb,    "Gradient Boosting")
\end{verbatim}

\subsection{Importância das variáveis (floresta aleatória)}

A importância relativa das variáveis é extraída diretamente do modelo e usada
para identificar os preditores mais relevantes do BVRP:

\begin{verbatim}
importances = pd.Series(
    rf.feature_importances_,
    index=FEATURES,
).sort_values(ascending=False)

importances.plot(kind="barh")
plt.xlabel("Importancia relativa")
plt.tight_layout()
\end{verbatim}

\section{Limpeza e controle de qualidade dos dados}

\begin{verbatim}
df = df.replace([np.inf, -np.inf], np.nan)
df = df.dropna(subset=["iv_30d", "rv_30d", "vrp_30d"])
df = df.sort_values("date").reset_index(drop=True)
\end{verbatim}

\section{Resumo}

Os trechos apresentados neste apêndice documentam os principais passos
computacionais da dissertação, permitindo replicar:

\begin{itemize}
    \item o cálculo da volatilidade realizada ($RV_{30}$) e do BVRP;
    \item a construção dos retornos futuros logarítmicos em cinco horizontes
          ($h \in \{1, 5, 10, 20, 30\}$ dias);
    \item a classificação das observações em regimes de baixa, média e alta
          volatilidade realizada ($q_{25} = 39{,}3\%$; $q_{75} = 65{,}6\%$);
    \item as regressões por MQO com erros-padrão de Newey--West e os modelos
          condicionais com interação BVRP $\times$ regime;
    \item a estimação dos quatro modelos supervisionados (LASSO, Ridge,
          \textit{floresta aleatória}, \textit{Gradient Boosting}) com avaliação
          genuína fora da amostra (divisão 70\,\%/30\,\%).
\end{itemize}

Os scripts completos estão disponíveis no repositório da dissertação e podem
ser estendidos para novos horizontes, conjuntos de variáveis ou algoritmos